\section{Filling-pair coefficients and a weighted mixing converse}
\label{sec:pairconverse}

The preceding converse criteria suggest retaining two independent
sets of curve data.  This section carries that proposal through three
concrete steps.  First, a tensor product of the coefficient construction
recovers marked filling pairs, hence square-tiled differentials.
Second, an exactly weighted count of Teichm\"uller orbit segments
unfolds into a time integral of correlations.  Third, an elementary
de-averaging argument converts an assumed counting remainder into
a mixing rate.  An arithmetic torus calculation illustrates the
additional cancellation needed to prove such a remainder.  The
identities and conditional implications below do not constitute a new
proof of Teichm\"uller mixing.

\subsection{Two independent character variables retain the pair}
\label{subsec:paircoefficients}

Keep a closed marked hyperbolic surface $X$ as an auxiliary device,
and let $\rho,\sigma$ be independent rank-two character variables.
The length-valued product of Section~\ref{sec:globalfilters} can be
tensored with itself, retaining two separate length coordinates:
\[
 \mathbf E_X(\rho)\boxtimes\mathbf E_X(\sigma)
 =\sum_{C,D} T_C(\rho)T_D(\sigma)\,
       (\nu_C^X\boxtimes\nu_D^X).
\]
Here $\boxtimes$ is the external product of locally finite discrete
measures, and both $C$ and $D$ range over integral weighted
multicurves.  The two variables are essential: multiplication in
one character algebra would resolve crossings between $C$ and $D$,
whereas their ordered pair is precisely the information required here.

\begin{proposition}[The first corner of a pair coefficient]
\label{prop:pairfirstcorner}
The coefficient of $T_C(\rho)T_D(\sigma)$ is supported in
\[
 [\ell_X(C),\infty)\times[\ell_X(D),\infty),
\]
and has an isolated atom of mass one at
$(\ell_X(C),\ell_X(D))$.  Its remaining support misses this corner.
\end{proposition}

\begin{proof}
The one-variable unit-diagonal identity
\eqref{eq:globaldiagonal} gives
\[
 \nu_C^X=\delta_{\ell_X(C)}+\tau_C^X,
 \qquad \operatorname{supp}\tau_C^X
        \subset(\ell_X(C),\infty),
\]
and similarly for $D$.  Expanding the external product proves the
claim; local finiteness isolates the corner.  Each finite rectangle
of source lengths involves only finitely many source configurations,
so the argument also applies before taking any transform.
\end{proof}

After this extraction, the marked labels retain $i(C,D)$ and the
topological filling condition.  Thus one can restrict to filling
pairs and apply weights to their associated flat differentials.
The auxiliary lengths $\ell_X$ provide the triangular extraction;
the area of the resulting differential will be $i(C,D)$, a different
quantity.  In particular, a bound on area is not a bound on both
auxiliary hyperbolic lengths on the marked space.  Passing to an
area cutoff and a mapping-class quotient is an additional operation,
whose quantitative cost is not supplied by fixed-coefficient
convergence.


\subsection{Filling pairs as actual coordinates for the flow}
\label{subsec:pairgeometry}

For definiteness let $S$ be closed, oriented, and of genus $g\geq2$.
A pair $(A,B)$ of nonzero measured foliations is \emph{filling} if
\[
 i(A,C)+i(B,C)>0\qquad\text{for every nonzero measured foliation }C.
\]
The condition is stronger than $i(A,B)>0$.  The realization theorem of
Gardiner--Masur, also proved by Wentworth using harmonic maps, identifies
filling pairs with marked nonzero holomorphic quadratic differentials
\cite[Theorem~1.6 and Corollary~1.7]{WentworthGardiner}; see also
\cite[Theorem~1.6]{AranaHerreraSquare2020} for the equivariant
homeomorphism and its integral interpretation.
Use the convention that $A=\mathcal V(q)$ is the vertical foliation
with transverse measure $|dx|$, and $B=\mathcal H(q)$ is the horizontal
foliation with transverse measure $|dy|$, in coordinates $q=dz^2$.
Write $q(A,B)$ for the resulting differential.  Then
\begin{align}
 \operatorname{Area}(q(A,B))&=i(A,B),\\
 g_tq(A,B)&=q(e^tA,e^{-t}B),
 \qquad g_t=\operatorname{diag}(e^t,e^{-t}).
 \label{eq:pairgeometryflow}
\end{align}
In particular
\[
 \widehat q(A,B):=
 q\left(\frac{A}{\sqrt{i(A,B)}},
         \frac{B}{\sqrt{i(A,B)}}\right)
\]
has unit area.  This is a genuine parametrization of points of the
flow space; the two foliations are the stable and unstable data of a
single differential, rather than two unrelated curves.

For an integral pair
$A=\sum_i a_i\alpha_i$, $B=\sum_j b_j\beta_j$ in minimal position,
there is a concrete realization.  Put a rectangle of horizontal width
$a_i$ and vertical height $b_j$ at each intersection of $\alpha_i$
with $\beta_j$, and glue the sides in the cyclic orders supplied by
the embedded curves.  The filling condition makes the complementary
regions disks.  Minimal position excludes bigons.  A complementary
$2k$-gon yields total angle $k\pi$, hence a zero of order $k-2$;
$k=2$ gives a regular point.  Integral side lengths subdivide the
rectangles into unit squares, with edge identifications by translations
or half-turns.  The total number of squares is
\[
 \sum_{i,j}a_i b_j i(\alpha_i,\beta_j)=i(A,B).
\]
Thus the discrete objects suggested by the pair coefficients are
precisely square-tiled half-translation surfaces.  Translation surfaces
form the subcollection in which the half-translation structure has
trivial linear holonomy.  Formula~\eqref{eq:pairgeometryflow} is then
visible by stretching and contracting the rectangles.

\subsection{The correct quotient and the two different counting scales}
\label{subsec:pairquotient}

Let $\Gamma=\operatorname{Mod}(S)$.  Equivariance of the pair
parametrization identifies the stabilizer of the measured pair with
the finite automorphism group of its differential:
\[
 \operatorname{Stab}_\Gamma(A,B)
 =\operatorname{Stab}_\Gamma(A)\cap\operatorname{Stab}_\Gamma(B).
\]
For integral pairs, finiteness also follows by acting on the finite
filling graph $A\cup B$ and then on its complementary disks; for
general measured foliations it follows from the realized conformal
automorphism group. It matters here that
one fixes the measured foliations, including their scales.  Stabilizers
of projective foliations can contain nontrivial flow translations.

One simultaneous $\Gamma$-orbit of a filling pair represents just one
point in moduli space and has constant intersection number.  To obtain
a growing family, prescribe the two curve types separately and take
the diagonal quotient:
\[
 \bigl(\Gamma\alpha\times\Gamma\beta\bigr)_{\mathrm{fill}}/\Gamma.
\]
If $\Gamma_\alpha=\operatorname{Stab}_\Gamma(\alpha)$ and
$\Gamma_\beta=\operatorname{Stab}_\Gamma(\beta)$, its classes are the
relevant double cosets in
$\Gamma_\alpha\backslash\Gamma/\Gamma_\beta$, represented by
$(\alpha,\phi\beta)$.  For an orbifold volume comparison the natural
weight of such a class is
\[
 \frac{1}{|\Gamma_\alpha\cap
                \phi\Gamma_\beta\phi^{-1}|}.
\]
This is the same stabilizer factor appearing in the square-tiled
count of \cite[Section~1 and Section~3]{AranaHerreraSquare2020}.

Fix a connected component $\mathcal Q$ of a stratum and put
$h=\dim_{\mathbb C}\mathcal Q$.  The real period-coordinate dimension
of its cone is $2h$, and the unit-area locus has real dimension $2h-1$.
If $D_r$ multiplies every period by $r$, then area is multiplied by
$r^2$.  With the appropriate period lattice and its covolume fixed,
Lebesgue measure $\nu$ scales by
\[
 \nu(D_r E)=r^{2h}\nu(E).
\]
Consequently an \emph{area} cutoff $L$ has natural count exponent $h$,
whereas a common \emph{linear} scale $R=\sqrt L$ has exponent $2h$.
These conventions must not be interchanged when converting powers
into rates in flow time.

For the closed principal quadratic stratum $h=6g-6$.  For a
quadratic stratum with $n$ specified singularities, including marked
regular points, $h=2g-2+n$ when the differential is not a square;
the corresponding square locus has $h=2g-1+n$.
Lower-dimensional strata require their own period-coordinate lattice
and Masur--Veech measure.  The ambient product of two full-dimensional
Thurston measures cannot serve as their volume form.  Likewise,
counting all square-tiled surfaces in a stratum does not by itself
approximate an arbitrary affine invariant measure on a smaller locus.

\subsection{The existing square-tiled counting theorem and the new target}
\label{subsec:pairliterature}

The passage from pairs to square tilings is established mathematics,
not a new consequence of the coefficient construction.
Arana-Herrera's effective theorem counts square-tiled surfaces in a
fixed connected component with prescribed horizontal and vertical
multicurve types, with a power-saving error in the area cutoff
\cite[Theorem~1.1]{AranaHerreraSquareEffective}.  Its proof explicitly
passes through Eskin--Mirzakhani--Mohammadi orbit counting and effective
horosphere equidistribution, using exponential Teichm\"uller mixing
\cite[Theorem~3.2 and its proof]{AranaHerreraSquareEffective}.
Using that estimate to recover mixing would therefore not provide an
independent proof of mixing.

There is a useful distinction already visible in that paper.
When one foliation type is fixed and the other is unrestricted,
standard lattice counting replaces the mapping-class-orbit estimate
\cite[Theorem~3.9]{AranaHerreraSquareEffective}.
This provides a tractable elementary benchmark.  It does not control
arbitrary initial and terminal observables at long times: their weights
must still be tracked through the reciprocal scaling in
\eqref{eq:pairgeometryflow}.


\subsection{What refining the coefficient mesh actually proves}

There are two independent limiting operations: approximation of a phase-space
integral by rational points, and factorization of that integral under a long
flow.  The first can be effective without implying the second.
Here is a precise local calculation.  Put
\[
 A_t=\operatorname{diag}(e^t I_d,e^{-t}I_d),\qquad
 H_t(x)=F(x)G(A_tx),
\]
where $F\in C_c^1(\mathbb R^{2d})$ has support in a fixed compact set $K$ and
$G\in C_b^1(\mathbb R^{2d})$.  For $R\geq1$ define
\[
 Q_R(t)=R^{-2d}\sum_{k\in\mathbb Z^{2d}}H_t(k/R),\qquad
 I(t)=\int_{\mathbb R^{2d}}H_t(x)\,dx.
\]
Then, for $t\geq0$,
\begin{equation}\label{eq:pairquadrature}
 |Q_R(t)-I(t)|\leq C_K R^{-1}e^t
       \|F\|_{C^1}\|G\|_{C^1}.
\end{equation}
Indeed, partition space into cubes of side $R^{-1}$, compare the value at one
corner with the integral on its cube, and use
\[
 \|\nabla H_t\|_\infty
 \leq \|\nabla F\|_\infty\|G\|_\infty
       +e^t\|F\|_\infty\|\nabla G\|_\infty.
\]
Only $O_K(R^{2d})$ cubes contribute.  A fixed smooth density can be absorbed
in $F$.  To model unit-area data by an ambient lattice, one can first use a
smooth radial cutoff in a compact area shell bounded away from area zero;
radial normalization is then smooth.  This local argument does not supply
uniform estimates across arbitrarily many period-coordinate charts, their
boundaries, or the cusp.

At the diagonal scale $R=e^t$, the displayed error has no decay.
Taking $R=e^{\sigma t}$, with $\sigma>1$, gives a quadrature error
$O(e^{-(\sigma-1)t})$.  This removes the resolution obstruction in this
particular local model.  It does not evaluate the remaining integral $I(t)$.
For Schwartz data, Poisson summation makes the distinction exact:
\begin{equation}\label{eq:pairpoissonzero}
 Q_R(t)=\sum_{n\in\mathbb Z^{2d}}\widehat H_t(Rn)
       =I(t)+\sum_{n\ne0}\widehat H_t(Rn).
\end{equation}
The zero Poisson frequency is the two-time integral itself, not the product
of its one-time means.  Suppressing the nonzero Poisson frequencies therefore
proves accurate quadrature, not mixing.  On the quotient, one must still
analyze the sum over changes of marking that returns $A_tK$ to a chosen
collection of charts.

Even exact sampling identities do not force mixing.  On
$\mathbb R/\mathbb Z$, with $g_t(x)=x+t$ and
$F(x)=G(x)=\cos(2\pi x)$,
\[
 \int_0^1F(x)G(g_tx)\,dx=\tfrac12\cos(2\pi t).
\]
For every integer $R\geq3$, the corresponding equal-mesh average is exactly
this integral, by finite Fourier orthogonality.  Both means vanish, while the
correlation does not decay.  This example isolates the logical issue; it is
not a model of the full geometry of a stratum.



\subsection{An exact weighted bisector kernel}
\label{subsec:twopointkernel}

There is a concrete discrete count whose remainder does imply mixing.
The necessary weight can be specified exactly, and the passage from a
cumulative count to a correlation has an elementary quantitative proof.
This construction uses the full unit cotangent bundle, with its
Masur--Veech measure supported on the principal stratum; it is not an
automatic statement for an arbitrary lower-dimensional stratum or
affine invariant measure.

Let $S$ be closed of genus $g\geq2$, put $h=6g-6$, and replace the
mapping class group by a torsion-free finite-index subgroup $\Gamma$.
This removes stabilizer factors. Write
\[
 Y=\Gamma\backslash Q^1\T(S),\qquad
 B=\Gamma\backslash\T(S),
\]
and normalize the invariant measure $\mu$ on $Y$ to have mass one.
Its projection $m$ to $B$ also has mass one. Use the same notation for
the lifted measures, and disintegrate
\[
 \dd\mu(q)=\dd s_x(q)\,\dd m(x),\qquad q\in S(x),
\]
where $S(x)=\pi^{-1}(x)$ and $s_x$ is a probability measure. The
Teichm\"uller polar-coordinate Jacobian $J_x(q,t)>0$ is defined almost
everywhere by
\begin{equation}
 \dd m\bigl(\pi(g_tq)\bigr)
   =J_x(q,t)\,\dd s_x(q)\,\dd t.
 \label{eq:twopointpolar}
\end{equation}
This is the polar-coordinate formula of
\cite[Section 2.4, equation (2.3)]{ABEMPolar}. We use the conventional
Teichm\"uller time $g_t=\diag(e^t,e^{-t})$.
The equality is an equality of measures under the map
$(q,t)\mapsto\pi(g_tq)$ on the marked universal space, with $t>0$;
it is not a claim that this map is one-to-one on the quotient $B$.

Let $D\subset\T(S)$ be a measurable fundamental domain. For $x,y\in D$
and $\gamma\in\Gamma$, let
\[
 \ell_\gamma=d_{\T}(x,\gamma y),\qquad
 q_\gamma=q_s(x,\gamma y)\in S(x)
\]
be the length and initial unit quadratic differential of the unique
Teichm\"uller segment. Its terminal differential, transported to
$S(y)$, is $\gamma^{-1}g_{\ell_\gamma}q_\gamma$.
For bounded observables $F,G$ on $Y$, understood through their
$\Gamma$-invariant lifts, define
\begin{equation}
 \begin{split}
 K_T^{F,G}(x,y)
  =\sum_{\substack{\gamma\in\Gamma\\1\leq\ell_\gamma\leq T}}
   \frac{e^{h\ell_\gamma}}{J_x(q_\gamma,\ell_\gamma)}\,
   F(q_\gamma)\,
   G\bigl(\gamma^{-1}g_{\ell_\gamma}q_\gamma\bigr).
 \end{split}
 \label{eq:twopointkernel}
\end{equation}
The lower cutoff avoids the polar singularity at zero. The sum is
finite for fixed $x,y,T$. Define it for $m\times m$-almost every
$(x,y)$ where the polar data are regular for all contributing
$\gamma$. Such a set has full measure, since $\Gamma$ is
countable; values on its complement can be assigned arbitrarily.
This convention makes no assertion about an exceptional prescribed
pair of basepoints.

\begin{proposition}[Exact unfolding of the weighted discrete count]
\label{prop:twopointunfolding}
For bounded $F,G$ and $T\geq1$,
\begin{equation}
 \begin{split}
 A_{F,G}(T)
 &:=\int_D\int_D K_T^{F,G}(x,y)\,\dd m(y)\,\dd m(x)\\
 &=\int_1^T e^{ht}
      \left(\int_Y F(q)G(g_tq)\,\dd\mu(q)\right)\dd t.
 \end{split}
 \label{eq:twopointexact}
\end{equation}
The same identity holds for $F,G\in L^2(\mu)$ whenever the
almost-everywhere kernel is interpreted through the corresponding
absolutely integrable unfolded expression.
\end{proposition}

\begin{proof}
First use nonnegative functions, so Tonelli's theorem applies.
The $y$-integral over $D$ and the sum over $\gamma$ unfold to an
integral over $\T(S)$. Substitute $z=\pi(g_tq)$ with $q\in S(x)$.
The factor $J_x(q,t)$ from~\eqref{eq:twopointpolar} cancels its
reciprocal in~\eqref{eq:twopointkernel}, giving
\[
 \int_D\int_1^T\int_{S(x)}
     e^{ht}F(q)G(g_tq)\,\dd s_x(q)\,\dd t\,\dd m(x).
\]
Disintegration gives~\eqref{eq:twopointexact}. The integral of the
absolute value is finite for bounded functions, and for $L^2$
functions it is at most
\[
 \|F\|_2\|G\|_2\int_1^T e^{ht}\,\dd t
\]
by invariance and Cauchy--Schwarz. Thus the signed or complex-valued
identity follows as well.
\end{proof}

The count is genuinely discrete before the basepoint integrals are
taken. It records both the initial and the terminal differential,
and its geometric weight is explicit. The identity does not prove
its asymptotic: that is the new counting theorem one would need.

\subsection{Removing the time average with a controlled loss}
\label{subsec:twopointdeaverage}

The next lemma applies to any probability-preserving flow. Write
$U_tG=G\circ g_t$ for its strongly continuous unitary action on
$L^2(\mu)$, and let $V$ be the infinitesimal generator, so that
$\partial_tU_tG=U_tVG$ on its domain. Put
\[
 C_{F,G}(t)=\int_Y F\,(G\circ g_t)\,\dd\mu,
 \qquad M_{F,G}=\mu(F)\mu(G).
\]
There is no complex conjugation in this notation; the usual
Hermitian correlation is obtained by replacing $F$ by $\overline F$.

\begin{theorem}[A cumulative counting converse]
\label{thm:twopointconverse}
Suppose, for some $\delta>0$, that
\begin{equation}
 \left|\int_1^T e^{ht}C_{F,G}(t)\,\dd t
       -M_{F,G}\frac{e^{hT}-e^h}{h}\right|
 \leq A(F,G)e^{(h-\delta)T},\qquad T\geq1.
 \label{eq:twopointcountassumption}
\end{equation}
Let $k\geq1$, $F\in L^2(\mu)$, and
$G\in\operatorname{Dom}(V^k)$. Then, for $t\geq2$,
\begin{equation}
 \begin{split}
 |C_{F,G}(t)-M_{F,G}|
 &\leq C_{h,\delta,k}
    \bigl(A(F,G)+\|F\|_2\|V^kG\|_2\bigr)
       e^{-\delta k t/(k+1)}.
 \end{split}
 \label{eq:twopointmixrate}
\end{equation}
Consequently a norm-controlled remainder for the discrete kernel
\eqref{eq:twopointkernel} gives a norm-controlled mixing estimate.
For one flow derivative the exponent is $\delta/2$. With arbitrarily
many controlled flow derivatives one obtains every exponent
$\eta<\delta$, at the cost of a higher regularity norm. No endpoint
rate $\delta$ follows from this single-test estimate alone;
Theorem~\ref{thm:twopointfullrate} below gives a stronger conclusion
for generator-differentiated endpoint weights.
\end{theorem}

\begin{proof}
Let $E(T)$ be the expression inside absolute values
in~\eqref{eq:twopointcountassumption}, and set
$R(T)=e^{-hT}E(T)$. Then $|R(T)|\leq A(F,G)e^{-\delta T}$ and,
distributionally (in fact continuously here),
\begin{equation}
 R'(T)+hR(T)=C_{F,G}(T)-M_{F,G}.
 \label{eq:twopointderivative}
\end{equation}
Choose a smooth, compactly supported kernel $\rho$ on $(-1,1)$ with
\[
 \int\rho(u)\,\dd u=1,\qquad
 \int u^j\rho(u)\,\dd u=0\quad(1\leq j<k),
\]
and put $\rho_\varepsilon(u)=\varepsilon^{-1}\rho(u/\varepsilon)$.
For $k>1$ this kernel need not be nonnegative. Integration by parts
in~\eqref{eq:twopointderivative} gives, for $0<\varepsilon\leq1$,
\[
 \left|\int\rho_\varepsilon(u)
       \bigl(C_{F,G}(t+u)-M_{F,G}\bigr)\,\dd u\right|
 \leq C_{h,\delta,k}A(F,G)\varepsilon^{-1}e^{-\delta t}.
\]
On the other hand,
\[
 C_{F,G}^{(k)}(t)
   =\int F\,U_tV^kG\,\dd\mu,
 \qquad |C_{F,G}^{(k)}(t)|
   \leq\|F\|_2\|V^kG\|_2.
\]
Taylor's formula and the vanishing moments therefore bound the
difference between the smoothed correlation and its value at $t$
by $C_k\varepsilon^k\|F\|_2\|V^kG\|_2$. Choosing
\[
 \varepsilon=e^{-\delta t/(k+1)}
\]
proves~\eqref{eq:twopointmixrate}.
\end{proof}

Only derivatives along the flow enter this last step. In particular,
it does not differentiate the spatially stretched observable
$G\circ g_t$ in an expanding transverse direction. This removes the
aliasing difficulty from the \emph{de-averaging step}; the difficulty
of proving the weighted count itself remains.

A convenient sufficient counting hypothesis is the almost-everywhere
bisector estimate
\begin{equation}
 \begin{split}
 K_T^{F,G}(x,y)
  &=\frac{e^{hT}-e^h}{h}\,
      f_F(x)f_G(y)+\mathcal E_T^{F,G}(x,y),\\
 f_F(x)&=\int_{S(x)}F\,\dd s_x,\qquad
 f_G(y)=\int_{S(y)}G\,\dd s_y,\\
 |\mathcal E_T^{F,G}(x,y)|
  &\leq C(x,y)\,\mathcal N(F)\mathcal N(G)
          e^{(h-\delta)T},\qquad
 C\in L^1(D\times D,m\times m).
 \end{split}
 \label{eq:twopointpointwise}
\end{equation}
Integrating gives~\eqref{eq:twopointcountassumption} with
$A(F,G)=\|C\|_1\mathcal N(F)\mathcal N(G)$. A direct estimate
for the averaged kernel is sufficient and is weaker than this
pointwise hypothesis. Uniformity over the observables, the two
basepoints, and all large real values of $T$ is part of the assertion;
one unweighted total count does not supply it.


A flow eigenfunction checks that the weighted hypothesis detects an
actual spectral obstruction. If an admissible nonzero test $G$ satisfies
$U_tG=e^{i\omega t}G$, $\mu(G)=0$, and $\|G\|_2=1$, take
$F=\overline G$. Then exact unfolding gives
\[
 A_{F,G}(T)
 =\frac{e^{(h+i\omega)T}-e^{h+i\omega}}{h+i\omega},
\]
whose size is comparable to $e^{hT}$ for large $T$. It violates
every positive-saving estimate in
\eqref{eq:twopointcountassumption}. This test is available only if
the asserted counting estimate covers the relevant observables;
an unweighted count does not test this eigenmode.

\subsection{A derivative endpoint weight preserves the full exponent}
\label{subsec:twopointderivativeweight}

For a counting estimate uniform over endpoint observables, there is
a stronger converse than temporal smoothing. One can differentiate
the \emph{endpoint test} inside the discrete count and use an exact
identity, without differentiating an asymptotic formula.

\begin{theorem}[A counting converse with no exponent loss]
\label{thm:twopointfullrate}
Use the notation of Proposition~\ref{prop:twopointunfolding}. Let
$F\in L^2(\mu)$ and $G\in\operatorname{Dom}(V)$, and suppose that
\[
 H_G=(h+V)G
\]
is an admissible terminal test for the weighted counting estimate.
Assume, for some $0<\delta\leq h$, that
\begin{equation}
 \left| A_{F,H_G}(T)
       -\mu(F)\mu(H_G)\frac{e^{hT}-e^h}{h}\right|
 \leq A(F,H_G)e^{(h-\delta)T},\qquad T\geq1.
 \label{eq:twopointderivativecount}
\end{equation}
Then the centered correlation
$\operatorname{Corr}_{F,G}(T)=C_{F,G}(T)-\mu(F)\mu(G)$ satisfies
\begin{equation}
 \begin{split}
 |\operatorname{Corr}_{F,G}(T)|
 &\leq A(F,(h+V)G)e^{-\delta T}\\
 &\quad+e^{-h(T-1)}
       |\operatorname{Corr}_{F,G}(1)|.
 \end{split}
 \label{eq:twopointfullrate}
\end{equation}
In particular, the counting saving $\delta$ gives the same
mixing exponent $\delta$.
\end{theorem}

\begin{proof}
Invariance of $\mu$ gives $\mu(VG)=0$, because the integral is a
continuous functional on $L^2(\mu)$ and
$\partial_tU_tG=U_tVG$. Hence $\mu(H_G)=h\mu(G)$.
The exact unfolding and the fundamental theorem of calculus give
\begin{align}
 A_{F,(h+V)G}(T)
  &=\int_1^T e^{ht}
       \bigl(hC_{F,G}(t)+C_{F,G}'(t)\bigr)\,\dd t\notag\\
  &=e^{hT}C_{F,G}(T)-e^h C_{F,G}(1).
 \label{eq:twopointfundamental}
\end{align}
Subtracting the main term in~\eqref{eq:twopointderivativecount}
therefore gives exactly
\begin{equation}
 E_{F,(h+V)G}(T)
  =e^{hT}\operatorname{Corr}_{F,G}(T)
       -e^h\operatorname{Corr}_{F,G}(1).
 \label{eq:twopointfullidentity}
\end{equation}
Divide by $e^{hT}$ and use the counting remainder.
\end{proof}

For example, if the count is uniform in a norm $\mathcal N$ with
\[
 A(F,H)\leq C\mathcal N(F)\mathcal N(H),
\]
the norm cost is explicit:
\begin{equation}
 \begin{split}
 |\operatorname{Corr}_{F,G}(T)|
 &\leq C\mathcal N(F)
       \bigl(h\mathcal N(G)+\mathcal N(VG)\bigr)e^{-\delta T}\\
 &\quad+e^{-h(T-1)}
       \|F-\mu(F)\|_2\|G-\mu(G)\|_2.
 \end{split}
 \label{eq:twopointfullnorm}
\end{equation}
For a suitable smooth norm on compactly supported observables,
this is one additional derivative along the flow in the terminal
test. On a noncompact space, any growth of the coefficients of
$V$ in the chosen norm must be retained in that estimate.
The terminal weight $(h+V)G$ can have either sign, so the counting
hypothesis must accommodate signed weights with a quantitative
norm bound.

The distinction from Theorem~\ref{thm:twopointconverse} is precise.
If only the primitive estimate for the particular pair $(F,G)$ is
available, temporal smoothing gives $\delta k/(k+1)$ with $k$
flow derivatives. If the same weighted counting theorem also
applies to $(F,(h+V)G)$, the exact identity above preserves
$\delta$. Thus exponent loss is not intrinsic to the converse
for a sufficiently rich family of endpoint weights.

In operator notation, the weighted time integral is governed by
$h+V$; inserting $(h+V)G$ removes this integral by an exact
boundary evaluation. No spectral gap is needed for this
algebraic step. The unresolved input is the norm-controlled
remainder for the corresponding discrete weighted count.

There is a corresponding improvement of the conditional cusp
calculation. Suppose bounded observables can be cut off at height
$H$ with correlation and mean errors $O(e^{-aH})$. Suppose the
counting bound for the cutoff pair
\[
 \bigl(F_H,(h+V)G_H\bigr)
\]
has constant $O(e^{bH})$, including the derivative of the cutoff.
Then~\eqref{eq:twopointfullrate} gives, after controlling the
boundary term,
\[
 |\operatorname{Corr}_{F,G}(t)|
   \ll e^{bH-\delta t}+e^{-aH}+e^{-ht}.
\]
For $a>0$, $b\geq0$, and $0<\delta\leq h$, choose
$H=\delta t/(a+b)$ to obtain
\begin{equation}
 \eta=\frac{\delta a}{a+b}.
 \label{eq:twopointfullcusprate}
\end{equation}
As before, unbounded observables require a norm-controlled tail
estimate in place of a bound on discarded measure alone.


\subsection{Comparison with the usual bisector count}
\label{subsec:twopointjacobian}

The exact weight $e^{ht}/J_x(q,t)$ cannot simply be omitted.
To compare normalizations, let $\mu_0$ be the local-product
normalization of~\cite[Section 2.2]{ABEMPolar} and put
$V_0=\mu_0(Y)$. Our probability normalization is
$\mu=\mu_0/V_0$. The conditional probability measures $s_x$ are
unchanged, whereas $m=m_0/V_0$ and $J_x=J_{0,x}/V_0$.
Use the canonical Hubbard--Masur factors $\lambda^\pm$ belonging
to $\mu_0$. For endpoints in a fixed compact subset of the principal
stratum, and for segments spending at least half their time in a
suitable larger compact set,
\cite[Proposition 2.5]{ABEMPolar} then gives
\begin{equation}
 J_x(q,t)
   =V_0^{-1}e^{ht}\lambda^-(q)\lambda^+(g_tq)
       +O(e^{(h-\alpha)t}).
 \label{eq:twopointHM}
\end{equation}
If the endpoint factors are bounded away from zero on the sets
under consideration, then, uniformly there for sufficiently large $t$,
the relative error in~\eqref{eq:twopointHM} is less than $1/2$.
Taking the reciprocal yields
\begin{equation}
 \frac{e^{ht}}{J_x(q,t)}
   =\frac{V_0}{\lambda^-(q)\lambda^+(g_tq)}
      +O(e^{-\alpha t}).
 \label{eq:twopointinverseHM}
\end{equation}
Thus the natural approximation is an ordinary bisector count with
the two endpoint observables divided by the corresponding
Hubbard--Masur factors, with the factor $V_0$. Usual bisector counts
often use the direction pointing back from the terminal point;
that convention replaces the terminal differential here by its
negative and requires the corresponding substitution in $G$.
There are three distinct estimates to check:
the weighted bisector remainder, the accumulated error in replacing
the exact Jacobian, and the contribution of segments failing the
recurrence or compactness conditions. Formula~\eqref{eq:twopointHM}
is not a global uniform asymptotic.

This optional comparison also has a dependency worth making
explicit: the proof of~\cite[Proposition 2.5 and Section 3.2]{ABEMPolar}
uses Forni's contraction estimate for the Hodge norm of the
Kontsevich--Zorich cocycle, quoted there as Theorem 3.3. The exact
Jacobian unfolding and the elementary de-averaging theorem above
do not use that estimate. Invoking the Hubbard--Masur asymptotic
would, however, retain this part of the existing cocycle machinery.

Effective bisector counts are already part of
\cite[Section 10]{AranaHerreraLattice}; those results are obtained
from exponential mixing. They locate the strength of the required
counting statement, but do not give an independent proof of mixing
when inserted into the converse above. An independent derivation
from coefficient distributions or a different counting mechanism
would have to establish the required weights and uniformity.

For completeness, the noncompact loss for the single-test smoothing
argument of Theorem~\ref{thm:twopointconverse} can also be computed. For bounded
observables, a discarded mass $O(e^{-aH})$ gives a correlation cutoff
error of the same order by invariance. For general $L^2$ observables
one must instead assume a corresponding norm-controlled tail estimate.
Suppose cutoffs at height $H$ have this error, the counting
constant for cutoff observables grows as $e^{bH}$, and their $k$th
flow-derivative bounds are uniform. That smoothing proof gives an
error of the form
\[
 \varepsilon^k+
      \varepsilon^{-1}e^{bH-\delta t}+e^{-aH}.
\]
Balancing the three terms yields
\begin{equation}
 \eta_k=\frac{\delta}{1+b/a+1/k}.
 \label{eq:twopointcusprate}
\end{equation}
This last formula is conditional on the stated cutoff estimates.
If the product of the relevant $k$th derivative norm and source
$L^2$ norm grows instead as $e^{dH}$, the same calculation gives
\[
 \eta_k=\frac{\delta}{1+b/a+(1+d/a)/k}.
\]
For fixed growth rates with $d=0$, increasing $k$ approaches
$\delta a/(a+b)$. These formulae do not establish any of the
required cusp or cutoff estimates.


\subsection{A completely explicit arithmetic two-endpoint calculation}


The modular case exhibits the additional cancellation that a two-endpoint
count must detect.  If
\[
 \gamma=\begin{pmatrix}a&b\\c&d\end{pmatrix}\in\mathrm{SL}_2(\mathbb Z),
 \qquad c>0,
\]
then
\begin{equation}\label{eq:pairbruhat}
 \gamma=
 \begin{pmatrix}1&a/c\\0&1\end{pmatrix}
 \begin{pmatrix}0&-1/c\\c&0\end{pmatrix}
 \begin{pmatrix}1&d/c\\0&1\end{pmatrix}.
\end{equation}
Modulo left and right integral unipotents, the two coordinates are
$a/c,d/c\pmod1$, with the coupled congruence $ad\equiv1\pmod c$.
They are initial and terminal horocyclic coordinates; the corresponding
cusp images are $\gamma\infty=a/c$ and
$\gamma^{-1}\infty=-d/c$.  They should not be identified, without an
additional unfolding argument, with arbitrary phase-space flow boxes.

Let $e(x)=e^{2\pi ix}$ and, for periodic test functions $f,g$ with
absolutely summable Fourier coefficients (in particular $C^2$ tests), set
\[
 D_c(f,g)=\frac1{\varphi(c)}
       \sum_{\substack{a\bmod c\\(a,c)=1}}
           f(a/c)g(\overline a/c),\qquad
 a\overline a\equiv1\pmod c.
\]
Writing $f(x)=\sum_m\widehat f(m)e(mx)$ and similarly for $g$ gives the exact
identity
\begin{equation}\label{eq:pairkloosterman}
 D_c(f,g)-\widehat f(0)\widehat g(0)
 =\frac1{\varphi(c)}
   \sum_{(m,n)\ne(0,0)}\widehat f(m)\widehat g(n)S(m,n;c),
\end{equation}
where
\[
 S(m,n;c)=\sum_{a\bmod c}^{*}e\bigl((ma+n\overline a)/c\bigr).
\]
The classical bound \cite[Section~1.1 and equation~(3.2)]{HumphriesModular}
\[
 |S(m,n;c)|\leq\tau(c)c^{1/2}(m,n,c)^{1/2}
\]
therefore supplies an actual product main term and a nontrivial remainder.
For example, defining
\[
 \mathcal A(f)=\sum_{m\in\mathbb Z}(1+|m|)^{1/2}|\widehat f(m)|,
\]
one obtains
\begin{equation}\label{eq:pairkloostermansmooth}
 |D_c(f,g)-(\int f)(\int g)|
 \leq\frac{\tau(c)c^{1/2}}{\varphi(c)}\mathcal A(f)\mathcal A(g)
 \ll_\epsilon c^{-1/2+\epsilon}\|f\|_{C^2}\|g\|_{C^2}.
\end{equation}
Here $(m,n)\ne(0,0)$ implies
$(m,n,c)^{1/2}\leq(1+|m|)^{1/2}(1+|n|)^{1/2}$, and two integrations by
parts bound $\mathcal A(f)$ by a constant times $\|f\|_{C^2}$.
For trigonometric polynomials supported on $|m|,|n|\leq M$, the sharper
explicit estimate is
\[
 |D_c(f,g)-(\int f)(\int g)|
 \leq\frac{\tau(c)\sqrt{cM}}{\varphi(c)}
             \|\widehat f\|_{\ell^1}\|\widehat g\|_{\ell^1},
 \qquad M\geq1.
\]
For a prime $c=p$ and $M<p$, the numerator $\tau(c)\sqrt{cM}$ can be
replaced by $2\sqrt p$; a frequency with exactly one zero coordinate gives
the Ramanujan sum $-1$.

Uniformity cannot be omitted: for $f_c(x)=g_c(x)=e(cx)$, both continuous
means vanish but $D_c(f_c,g_c)=1$.  The growth of the test norms in
\eqref{eq:pairkloostermansmooth} correctly prevents a false conclusion.
This is the arithmetic counterpart of the expanding-frequency problem.
The result is a rigorous benchmark for endpoint-residue decorrelation.  A
proof of mixing of the modular geodesic flow still requires the appropriate
unfolding, radial/time weights, and cusp estimates; no higher-genus analogue
of the displayed Kloosterman bound has been established here.

For comparison, fix two primitive oriented torus curve types.
In the corresponding matrix model, with positive oriented intersection
$c$, the remaining residue is $a\pmod c$, $(a,c)=1$, so there are exactly
$\varphi(c)$ possibilities.  Thus the unweighted count already has the
power-saving asymptotic
\[
 \sum_{c\leq L}\varphi(c)
   =\frac{3}{\pi^2}L^2+O(L\log L).
\]
Indeed, substitute $\varphi(c)=c\sum_{d\mid c}\mu(d)/d$ and sum first over
$c=dk$; the leading term is
$\tfrac12L^2\sum_{d\leq L}\mu(d)/d^2$
and the remainder is $O(L\log L)$.
These constants belong to this oriented matrix model; no unmentioned
quotient by signs or orbifold automorphism weight is being included.
The exact formula $\varphi(c)$ by itself says nothing about the values of
$S(m,n;c)$ for $(m,n)\ne(0,0)$.


\subsection*{The remaining independent estimate}

We have obtained an exact encoding of the filling-pair data, an exact
weighted orbit-counting identity, and a quantitative implication from
a remainder for that count to mixing.  What has not been established
is a norm-controlled estimate of the form
\eqref{eq:twopointcountassumption} from the coefficient distributions.
The tensor product and the first-corner extraction do not by themselves
estimate a sum over returning mapping classes.
In particular, an exact or quantitative comparison between the resulting
square-tiled sums and the bisector kernel has not yet been supplied.

The next analytic target can now be stated without referring vaguely
to a good scalar counting error: construct that comparison and control
the nonconstant two-endpoint weights in \eqref{eq:twopointkernel}, with the exact inverse Jacobian
or with a proved replacement and a controlled discarded contribution.
The modular calculation shows what success looks like.  Its scalar
count determines the zero mode, while the Kloosterman estimate controls
the coupled nonzero modes. No analogous estimate for the principal
higher-genus kernel is proved here. Section~\ref{sec:genus2markedtorus}
does close a different, explicitly defined return calculation on the
exceptional marked-torus model of the nonsquare genus-two $(2,2)$ stratum.

An independent proof of the displayed remainder would yield a new
route to mixing for the specified Masur--Veech measure, with the full rate
in \eqref{eq:twopointfullrate} when derivative endpoint weights are
controlled, and the cutoff rate~\eqref{eq:twopointfullcusprate}.
The single-test estimate~\eqref{eq:twopointmixrate} remains an
alternative under its weaker hypothesis.  It would not by itself classify invariant
measures or orbit closures, nor would it establish spectral properties
of a separate cocycle.  Those require additional information beyond
correlations of the base flow.


